\chapter{Discussion and Future Improvements}\label{ch:discussion}

\section{Strengths}

\begin{itemize}
  \item \textbf{Measured, not claimed.} Every gain comes from the same data and workload on
        both designs, with plans, raw JSON and test outputs in the repository.
  \item \textbf{Retention as metadata.} Archive and purge no longer depend on data volume. Their
        cost stays flat as the business grows, while the baseline grows linearly.
  \item \textbf{Safety by design.} Guards refuse illegal purges, frozen data is read-only, the
        application can only write the hot window, every operation is logged and resumable,
        and integrity is proved against control totals. The job stops if a frozen month
        changes.
  \item \textbf{Operations.} One command per month, a dry-run mode, the policy stored as data,
        and readable partition names.
\end{itemize}

\section{Limitations and remaining risks}

\begin{itemize}
  \item \textbf{Licence.} Partitioning is a paid Enterprise Edition option. On Standard
        Edition~2 the design would fall back to one table per month behind a \code{UNION ALL}
        view, without basic compression (EE only) and with weaker pruning.
  \item \textbf{Queries without a date filter} gain nothing (Q6), and multi-month lookups
        through local indexes read a few more blocks (Q2).
  \item \textbf{Warm cost.} It is probably dominated by maintenance of the global primary key (not measured separately); the
        monthly job is 8.9$\times$ faster, not the 10$\times$ we targeted.
  \item \textbf{No \code{record\_id} index on archived rows}: a lookup by \code{record\_id}
        through \code{BILLING\_ALL} scans the archive partitions.
  \item \textbf{Concurrency.} The application can still hold locks (its \code{INSERT} grant
        allows \code{LOCK TABLE}), and two lifecycle runs are not serialised. Guards and row
        counts prevent data loss, but a run can fail and must be re-run.
  \item \textbf{Evaluation context.} One laptop, single-user measurements, synthetic data
        (uniform customers, no corrections of old bills) and 24~million rows. Absolute times
        are indicative; logical reads and redo are the reliable comparisons.
  \item \textbf{Owner privileges.} \code{TELCO\_OWNER} can change the policy and read
        PDB-wide \code{V\$SQL}. In production it would be locked, and the job would run under
        a scheduler with auditing on \code{DLM\_CONFIG}.
\end{itemize}

\section{Possible improvements}\label{sec:improvements}

\begin{enumerate}
  \item \textbf{Reduce the warm cost:} use \code{NOLOGGING} moves followed by a backup (only
        without a standby), warm several months in parallel off-peak, or switch to a local
        primary key on \code{(record\_id, bill\_date)} if the application can accept it.
  \item \textbf{Schedule and serialise} the monthly job with \code{DBMS\_SCHEDULER} and a
        \code{DBMS\_LOCK} lock, with alerting on \code{ERROR} or \code{REFUSED} log entries.
  \item \textbf{Archive statistics:} gather table-level synopses on \code{BILLING\_XCHG}
        before each exchange, so exchanged partitions arrive with incremental synopses.
  \item \textbf{Tiered storage:} place \code{TELCO\_ARCH} on cheaper storage and make archive
        tablespaces read-only, so backups of history are taken once.
  \item \textbf{Measure under concurrency:} run the application workload during the lifecycle
        job, and test on production-scale volumes.
  \item \textbf{Audit:} Unified Auditing policies on \code{DLM\_CONFIG} and on lifecycle
        procedures.
\end{enumerate}
